\section{Related Work}
\label{sec:related}

Hash and sort-merge joins have a long systems history; Graefe surveys the
implementation trade-offs that remain relevant to partitioning and spill
~\cite{graefe1993query}.  Main-memory studies revisit those trade-offs under
modern processors and parallelism~\cite{blanas2011design,balkesen2013main}.
\system{} builds on partitioned hashing but changes routing at partition
granularity during execution.

Adaptive query processing uses observations to revise execution in the
presence of uncertain statistics~\cite{deshpande2007adaptive}.  Eddies route
tuples adaptively among operators~\cite{avni2000eddies}, and later systems
interleave planning with execution~\cite{trummer2019skinner}.  Our focus is
complementary: a fixed binary join with a bounded buffer, explicit distributed
migration, and skew-aware routing.

Parallel database engines use exchange operators, work stealing, and fine-grain
tasks to balance CPU work.  Morsel-driven execution demonstrates the value of
NUMA-aware scheduling and small work units~\cite{leis2014morsel}.  Vectorized
execution amortizes interpretation and function-call overhead
~\cite{boncz2005monetdb}.  \system{} assumes these mechanisms and adds a
control plane for deciding where join work should run.

Distributed streaming systems provide epoch, barrier, and backpressure ideas
for stateful dataflows.  \system{} adapts those ideas to a finite relational
operator: stable source offsets and epoch visibility make mode changes
recoverable without turning every tuple into a durable log record.
